% UPPER PRELIMINARIES BEGIN
\subsection{The shuffle graph}
Use the path graph and componentwise sorting representation of
Lemma~\ref{lem:matrix-graph}. Its components are the nonempty
equal-word intervals, including isolated vertices. For a fixed path
graph $G$, denote the componentwise-sorted uniform law by $P_G$.
Thus $P_N^K=\E P_G$.

% Lean: GraphOverlap/Density.lean: graphLaw_density, overlapIdentity;
% GraphOverlap/RunProduct.lean: boundedComponentOverlap_strong, boundedComponentOverlap.
\begin{lemma}\label{lem:upper-overlap}
The density of $P_G$ with respect to $\nuUnif$ is
\[
 f_G(\sigma)=\mathbf1_{\{\sigma\text{ is increasing on every component of }G\}}
                   \prod_{I\in\operatorname{comp}(G)}|I|!.
\]
For any two path graphs $G,H$,
\begin{equation}\label{eq:upper-overlap}
 \E_{\nuUnif}f_G f_H
 =\frac{\displaystyle\prod_{I\in\operatorname{comp}(G)}|I|!
                    \prod_{J\in\operatorname{comp}(H)}|J|!}
          {\displaystyle\prod_{I\in\operatorname{comp}(G\cup H)}|I|!}.
\end{equation}
If the components of $H$ have at most $B$ vertices, this is at most
$B^{|G\cap H|}$. Consequently, for $B\ge2$, it is also at most
$\binom B{\lfloor B/2\rfloor}^{|G\cap H|}$.
\end{lemma}
\begin{proof}
A permutation increasing on every component has exactly
$\prod_I|I|!$ preimages under componentwise sorting; a permutation
which is not increasing has none. This proves the density formula.
The simultaneous increasing condition for $G$ and $H$ is the increasing
condition for $G\cup H$, whose uniform probability is
$1/\prod_{I\in\operatorname{comp}(G\cup H)}|I|!$. This proves
\eqref{eq:upper-overlap}.

For a vertex $i$, let $r_G(i)$ be its position, counted from the left,
in its component of $G$. Then $\prod_i r_G(i)=\prod_I|I|!$ and
$r_G(i),r_H(i)\le r_{G\cup H}(i)$. Unless the edge immediately to
the left of $i$ belongs to both graphs, at least one of $r_G(i)$ and
$r_H(i)$ is one. At the other vertices, $r_H(i)\le B$. Thus
\[
 r_G(i)r_H(i)\le r_{G\cup H}(i)
       B^{\mathbf1_{\{(i-1,i)\in G\cap H\}}}.
\]
Multiplying over the vertices proves the claimed bound. Finally
$B=\binom B1\le\binom B{\lfloor B/2\rfloor}$.
\end{proof}

% Lean: GraphOverlap/Truncation.lean: truncatedTV, truncatedVariance_eq.
\begin{lemma}\label{lem:upper-truncation}
Let $S$ be any event determined by one word matrix, and let $G,G'$ be
the graphs of two independent matrices with the same deterministic
column counts. Then
\begin{equation}\label{eq:upper-truncation}
 \dTV(P_N^K,\nuUnif)
 \le\frac12\sqrt{\E\left[
   \bigl(\E_{\nuUnif}f_Gf_{G'}-1\bigr)
                      \mathbf1_S\mathbf1_{S'}\right]}+\Pp(S^c).
\end{equation}
The expression under the square root is nonnegative.
\end{lemma}
\begin{proof}
Write $h=\E[f_G\mathbf1_S]$ and $r=\E[f_G\mathbf1_{S^c}]$,
so the shuffle density is $h+r$. Since $f_G$ has mean one,
$\E_{\nuUnif}h=\Pp(S)$ and $\E_{\nuUnif}r=\Pp(S^c)$.
The triangle inequality gives
\[
 \E_{\nuUnif}|h+r-1|
 \le\E_{\nuUnif}|h-\Pp(S)|+2\Pp(S^c).
\]
Apply Cauchy--Schwarz to the first term. Independence and finite
summation identify its squared upper bound as
\[
 \E_{\nuUnif}(h-\Pp(S))^2
 =\E[(\E_{\nuUnif}f_Gf_{G'}-1)\mathbf1_S\mathbf1_{S'}].
\]
This also proves the nonnegativity assertion.
\end{proof}

% Lean: GraphOverlap/FactorialTail.lean; GraphOverlap/Corollaries.lean:
% boundedComponentTail_factorial, exists_boundedComponentTail.
\begin{lemma}\label{lem:upper-components}
Put $\mu_*=N\prod_{j=1}^K\max_d p_d^{(j)}$.
For every integer $B\ge1$,
\begin{equation}\label{eq:upper-components}
 \Pp(G\text{ has a component with more than }B\text{ vertices})
       \le\frac{N\mu_*^B}{(B+1)!}.
\end{equation}
In particular, if $\mu_*\le N^{-\eta}$ and $\eta B\ge A+1$,
the right side is at most $N^{-A}/(B+1)!$.
\end{lemma}
\begin{proof}
A component with more than $B$ vertices supplies $B+1$ distinct rows
having the same full word. The event is impossible if $B+1>N$,
so suppose $B+1\le N$.
For any specified set of $B+1$ rows, its probability
of agreeing in column $j$ is
\[
 \sum_d\frac{(Np_d^{(j)})_{B+1}}{(N)_{B+1}}
 \le\sum_d(p_d^{(j)})^{B+1}
 \le(\max_d p_d^{(j)})^B.
\]
Here $(m)_r=m(m-1)\cdots(m-r+1)$, and the first inequality follows
factor by factor; impossible samples have probability zero. Multiply
over columns and sum over the $\binom N{B+1}$ sets of rows.
Since $\binom N{B+1}\le N^{B+1}/(B+1)!$, this gives the result.
\end{proof}
% UPPER PRELIMINARIES END

% UPPER MAIN BEGIN
\section{Upper bound}\label{sec:upper-bound}\label{sec:upper}
\subsection{Empirical costs and an initial good shuffle}
We first give an upper bound requiring no limiting time profile.
The original deck size remains the logarithmic scale when the argument
is applied to a smaller population. Thus let $2\le n\le N$ and let
the current columns be indexed by $a<j\le K$, with $a<K\le T\log N$.
Their proportions, relative to $n$, are denoted by $p^{(j)}$. At integer
endpoints set
\begin{align*}
 A^{(N)}\left(\frac{s}{\log N},\frac{u}{\log N}\right)
       &=\frac1{\log N}\sum_{j=s+1}^u a_{p^{(j)}},\\
 B^{(N)}_\theta\left(\frac{s}{\log N},\frac{u}{\log N}\right)
       &=\frac1{\log N}\sum_{j=s+1}^u\psi_{p^{(j)}}(\theta),\\
 \beta^{(N)}\left(\frac{s}{\log N}\right)
 &=\frac{\log n}{\log N}
       -A^{(N)}\left(\frac{a}{\log N},\frac{s}{\log N}\right).
\end{align*}
We say that the empirical margin is at least $\eta$ if the following
inequalities hold for every integer pair $a\le s\le u\le K$:
\begin{equation}\label{eq:upper-empirical-margin}
 \begin{split}
 \beta^{(N)}\left(\frac K{\log N}\right)&\le-\eta,\\
 \inf_{3\le\theta\le4}\left\{
       \frac{\theta+3}{2}\beta^{(N)}\left(\frac s{\log N}\right)
       -B^{(N)}_\theta\left(\frac s{\log N},\frac u{\log N}\right)
       -2B^{(N)}_2\left(\frac u{\log N},\frac K{\log N}\right)\right\}
       &\le-\eta.
 \end{split}
\end{equation}
Only grid endpoints are required in this condition; the interpolation
between them is not an additional hypothesis. Zero pile proportions
are allowed. A column is $g$-good if two distinct coordinates of its
proportion vector are at least $g$.

% Lean: EmpiricalUpper/Restart.lean: terminal_horizon, restartMargin,
% restartHypotheses; Exploration/WindowClusters.lean: integerWindow_cluster.
\begin{lemma}\label{lem:upper-restart}
Fix $k\ge2$, $\eta_0,\eta,g,\gamma,\rho>0$, with
$4\eta\le\eta_0$, $\gamma\le1$ and
$\rho\log k\le\min\{\eta/7,1/4\}$.
For all sufficiently large $N$, an array on $N$ labels satisfying
\eqref{eq:upper-empirical-margin} with margin $\eta_0$ and the local
$(g,\gamma,\rho)$-good condition on $[0,K/\log N]$ has a suffix
with the following properties. Its first column is $g$-good; two more
$g$-good columns occur within the next $\lceil6/\gamma\rceil$ columns;
its empirical margin, after resetting its initial index to zero, is
at least $2\eta$; and every consecutive window of
$\lceil\rho\log N\rceil+1$ columns contains at least
$\gamma\rho\log N-2$ good columns. Its maximal full-word mean is at
most $N^{-2\eta}$. The total variation distance for the original array
is no larger than that for this suffix.
\end{lemma}
\begin{proof}
Since $a_p\le\log k$, the terminal margin implies
$(1+\eta_0)\log N\le K\log k$. The initial window of
$\lceil\rho\log N\rceil+1$ columns therefore fits for large $N$.
Passing from a closed time window to the integer window loses at most
two columns. Among at least $\gamma\rho\log N-2$ good indices in
that window there are three columns, say $r+1<j+1<h+1$, with
$h-r\le\lceil6/\gamma\rceil$. Indeed, otherwise every third good
index would be separated by more than that amount, contradicting
the number of good indices once $\gamma\rho\log N>8$.
Delete the $r$ columns preceding column $r+1$.

Deleting $r$ columns increases $\beta^{(N)}(s/\log N)$ at subsequent times
by at most $r\log k/\log N$, and increases every expression in
\eqref{eq:upper-empirical-margin} by at most
$7r\log k/(2\log N)$. Here $r\le\rho\log N+2$; hence the total
loss is at most $\eta/2+7\log k/\log N\le\eta$ for large $N$.
The retained margin is thus at least $2\eta$. Every retained integer
window is an original integer window, so its good-column bound is
inherited. The identity
$\mu_*=N^{\beta^{(N)}(K/\log N)}$ gives the asserted maximal mean.
Finally, the deleted shuffles act by averaging translates of the
retained permutation law. Uniform measure is invariant under these
translations, and total variation contracts under averaging. This
proves the last assertion.
\end{proof}

% Lean: Contracts/V1/Empirical.lean: E; EmpiricalUpper/Main.lean: empiricalUpper.

\begin{proof}[Proof of Theorem~\ref{thm:upper-empirical}]
Put $\eta=\min\{\eta_0/4,1\}$,
$g=\min\{g_0,1/(2k)\}$ and $\gamma=\min\{\gamma_0,1\}$.
The local condition remains true with these smaller parameters.
Choose $\rho_0$ small enough for Proposition~\ref{prop:upper-conditional}
and so that $\rho_0\log k\le\min\{\eta/7,3\eta/160,1/4\}$.
Lemma~\ref{lem:upper-restart} gives a suffix with margin $2\eta$,
the required good cluster and integer-window condition, and maximal
mean at most $N^{-2\eta}$.

Proposition~\ref{prop:upper-conditional} and the construction in
Section~\ref{sec:upper-exponential} supply $2k$ adapted explorations,
each with conditional expected remainder at most $CN^{-3\eta/80}$.
The single-matrix regularity event used there satisfies
\[
 \Pp(\mathcal R^c)\le
 AN^A e^{-c_H N^{\eta/30}}
 +2N\sqrt{AN^A e^{-c_H N^{\eta/30}}}
 +\frac{k}{2}N^{\eta/10-2\eta}.
\]
Choose an integer $B\ge2$ with $\eta B\ge2$ and intersect
$\mathcal R$ with the event that all graph components have size at
most $B$. Lemma~\ref{lem:upper-components} adds at most
$N^{-1}/(B+1)!$ to the failure probability. Apply
Proposition~\ref{prop:upper-exponential} with $d=2k$ and
$b=\binom B{\lfloor B/2\rfloor}$. All its hypotheses concern the
actual pair histories constructed below, and its denominator tends
to one. The resulting bound is a fixed multiple of
$N^{-3\eta/160}+N^{-19\eta/10}+N^{-1}$, plus stretched exponential
terms. After decreasing the positive exponent and increasing $N_0$,
this is at most $N^{-c}$. The contraction in
Lemma~\ref{lem:upper-restart} completes the proof.
\end{proof}

\subsection{A limiting profile}
% Lean: ProfileUpper/Buffered.lean: bufferedOfEmpirical;
% ProfileUpper/Main.lean: bufferedUpper; Margins/ProfileMargin.lean.

\begin{proof}[Proof of Theorem~\ref{thm:upper-profile}]
Use Corollary~\ref{cor:fixed-partition-transfer} with accuracy
$\eta/100$, first fixing $\chi$, then $\rho$, and finally $\varphi$.
The comparison is uniform over $\theta\in[2,4]$ and over interval
endpoints. Restricting to grid endpoints is exact; changing the
horizon from $T$ to $\lfloor T\log N\rfloor/\log N$ costs at most
$14\log k/\log N$. Consequently the two profile inequalities
give \eqref{eq:upper-empirical-margin} with margin $\eta/2$ for large
$N$. Choose the upper bound on $\rho$ also to satisfy
Theorem~\ref{thm:upper-empirical} with this margin and the given
$g,\gamma$. Since $K/\log N\le T$, the local condition on $[0,T]$
implies the one on $[0,K/\log N]$. The empirical theorem gives
$N^{-c}$, which is at most $\varepsilon$ for sufficiently large $N$.
All thresholds are uniform over the arrays in the statement.
\end{proof}

% Lean: Shared/Profile.lean: ArbitrarilyFineApproximation, criticalTime;
% ProfileUpper/Critical.lean: criticalOfBuffered, marginMonotonicity (Analytic.lean).

\begin{proof}[Proof of Corollary~\ref{cor:upper-critical}]
If $t$ exceeds the displayed infimum, its definition supplies $0<s<t$ at which
the reduced margin is negative. The reduced margin is nonincreasing
in the horizon, so it is negative at $t$ as well. Apply
Theorem~\ref{thm:upper-profile}, selecting the partition beneath its
mesh bound, then $\rho$ beneath both allowed bounds, and finally an
allowed positive $\varphi$. Arbitrarily fine approximation supplies
the eventual deterministic hypotheses, proving convergence.
For every fixed $t$ exceeding the infimum, convergence implies eventually
$t_{\rm mix}^{(N)}(\varepsilon)\le\lfloor t\log N\rfloor$.
Divide by $\log N$, take the upper limit, and take the infimum over
these $t$. If the infimum is $+\infty$, the upper-limit assertion
holds in the extended order.
\end{proof}

\subsection{Random pile sizes}
% Lean: Contracts/V1/Profile.lean: REmpiricalBound, REmpirical, RProfile;
% ProfileUpper/Mixture.lean, ProfileUpper/Random.lean, ProfileUpper/Main.lean: randomUpper.
\begin{corollary}\label{cor:upper-random}
Under the parameter choices of Theorem~\ref{thm:upper-empirical}, let
$Q_N$ be any law on deterministic arrays of length $K\le T\log N$.
For all sufficiently large $N$,
\[
 \dTV\left(\int P_N^K(\mathbf n)\,Q_N(d\mathbf n),
                         \nuUnif\right)
 \le Q_N\{\text{the empirical margin or local-good condition fails}\}
                        +N^{-c}.
\]
Thus the annealed distance tends to zero if these two conditions hold
with probability tending to one, with one common deterministic $\rho$.
The same conclusion holds under the parameter choices of
Theorem~\ref{thm:upper-profile} if almost-likeness and local goodness
each hold with probability tending to one.
\end{corollary}
\begin{proof}
Convexity of total variation bounds the mixture distance by the
$Q_N$-average of the conditional distances. On the good arrays use
$N^{-c}$, and on the other arrays use the bound one. This proves the
finite inequality and its empirical consequence. In the profile
case, for any fixed positive error use
Theorem~\ref{thm:upper-profile} on the intersection of the two good
events. Its complement has probability tending to zero by the union
bound. Letting the fixed error tend to zero proves the claim.
\end{proof}
% UPPER MAIN END

% UPPER FIRST MOMENT BEGIN
\section{The first moment}\label{sec:upper-first}
We prove the uniform estimate used for the whole deck and, later, for
the random populations left inside unexplored cylinders.

% Lean: Contracts/V1/Empirical.lean: F; FirstMoment/Empirical.lean: empiricalFirstMoment;
% FirstMoment/UniformGraph.lean: empiricalGraphDecay.
\begin{proposition}\label{prop:upper-first}
Fix $k\ge2$, $T,\eta,g>0$, with $g<1/k$. There is $N_0\ge2$ such
that, for every $N\ge N_0$, $2\le n\le N$, and
$a<K\le T\log N$, the following holds. Suppose that the first of
the fixed-count columns on $n$ labels is $g$-good, and that
\eqref{eq:upper-empirical-margin} holds. For independent shuffle
graphs $G,G'$ of these columns,
\begin{equation}\label{eq:upper-first}
 \E|G\cap G'|\le N^{-\eta/8}.
\end{equation}
\end{proposition}

\subsection{Interval fluctuations}
Condition on the assignment of the first column. A nonempty pile
$R$ has size $M\le(1-g)n$. The remaining columns are still uniform
fixed-count columns on all $n$ labels. The labels outside $R$ will
provide the unused population needed to smooth interval endpoints.
Permutation invariance of the rows shows that the graph law within
$R$ depends only on $M$ and its deterministic rank offset.

For a legal prefix $x$ following this first digit, write $\mu_x$ for
its mean count in $R$, $\ell_x=\log_N\mu_x$, and
$e_x=-B^{(N)}_2(M_x/\log N,K/\log N)$, where $M_x$ is its absolute final column.
Its descendants occupy an interval with endpoints $\iota_x,\tau_x$,
and its internal equal-word edges form a graph $G_x$.
All counts and fluctuations in this paragraph exclude the already
fixed first-column assignment.

Fix $0<\chi<1/k$, and in each subsequent column let $l_j,r_j$ be the
smallest and largest digits of mass at least $\chi$. For a full word,
its left boundary stops at its first digit strictly greater than $l_j$,
and its right boundary at its first digit strictly less than $r_j$;
if no crossing occurs, use the full word. Let $v$ be the longer
boundary. If either boundary has full length, call the pair terminal;
its block is the single full word. Otherwise the pair has two genuine
crossings. For every column through $v$, consider the rows whose first
difference from $v$ occurs there and is smaller, respectively larger.
Let $L_v,R_v$ be the largest means of these pieces, and put
\[
 F_v=\tfrac12\log_N\min\{L_v,R_v\},\qquad V_v=\ell_v.
\]
Both side means are positive. Their elementary comparison is
\begin{equation}\label{eq:upper-side-scale}
 \min\{L_v,R_v\}\ge\chi\mu_v,\qquad
 M\le(2|v|+\chi^{-1})\max\{L_v,R_v\}.
\end{equation}
Indeed, each genuine crossing contributes a side piece of mean at
least $\chi\mu_v$. The at most $2|v|$ first-difference pieces together
with the prefix cylinder partition $R$ in expectation, giving the
second inequality. In particular,
$2F_v\ge V_v-\log(1/\chi)/\log N$.

% Lean: Hypergeometric/AntiConcentration.lean: pointMass_le_four_div_sqrt;
% Hypergeometric/Cylinder.lean: anchorCylinderTail;
% lower_bound_v2/Thinning/Tails.lean: bernstein_of_transform, half_mean_of_transform.
\begin{lemma}\label{lem:upper-atoms}
If $Y$ is hypergeometric with population $n$, marked count $r$, and
sample size $m$, then
\[
 \E Y=\frac{rm}{n},\qquad
 \Var(Y)=\frac{rm(n-r)(n-m)}{n^2(n-1)}\quad(n\ge2).
\]
If $\sigma^2=\Var(Y)>0$, then $\sup_z\Pp(Y=z)\le4/\sigma$.
A prefix count in any deterministic set of rows, with mean $\mu$,
satisfies, for $t>0$,
\[
 \Pp(|X-\mu|\ge t)\le2\exp\left\{-\frac{t^2}{4(\mu+t)}\right\},
 \qquad
\Pp(X\le\mu/2)\le\exp\left\{-\frac{1-\log2}{2}\mu\right\}.
\]
\end{lemma}
\begin{proof}
The moment formulas follow by expanding a sum of marked-position
indicators and using the one- and two-position sampling probabilities.
For the atom estimate let $p_q=\Pp(Y=q)$ and define, within this proof,
$b(x)=(r-x)(m-x)/n$ and
$d(x)=x(n-r-m+x)/n$. The binomial coefficient formula gives detailed
balance $b(q)p_q=d(q+1)p_{q+1}$, while
$d(x)-b(x)=x-\E Y$. Telescoping therefore yields
\[
 b(q)p_q=\sum_{j\le q}(\E Y-j)p_j
       \le\E|Y-\E Y|\le\sigma.
\]
Choose a mode $q$. Comparing its mass with the preceding mass by
detailed balance gives $q\le\E Y+2$. Since
$0\le r+m-q-\E Y\le2n$, expanding the quadratic $b$ gives
$b(q)\ge b(\E Y)-4$. Also
$b(\E Y)=\sigma^2(n-1)/n\ge\sigma^2/2$ when $n\ge2$.
If $\sigma^2\ge16$, then $b(q)\ge\sigma^2/4$, and the displayed
inequality gives $p_q\le4/\sigma$. If $\sigma^2<16$, use
$p_q\le1\le4/\sigma$.

For the prefix bound, expose successive independent uniform subsets
corresponding to its digits. The falling-factorial expansion of the
exponential transform gives the sampling bound for a nonnegative
exponent. Apply the same expansion to the complement sample for a
negative exponent, and then iterate the hypergeometric thinning.
This gives $\E e^{tX}\le\exp\{\mu(e^t-1)\}$ for every real $t$.
The centered transform and $e^t-1-t\le t^2$ for $|t|\le1$,
with $|t|$ chosen as the deviation divided by $2(\mu+\text{deviation})$,
give the two-sided bound. For the lower half-mean bound use
$t=-\log2$ in the uncentered transform. An arbitrary deterministic
row set has the same count law as a uniform row set of that size by
permutation invariance, so the argument applies to the anchor $R$.
\end{proof}

% Lean: Anchor/Smoothing.lean: endpointSmoothing, endpointShiftHypergeometric,
% smoothingFailure, endpointVariance; Anchor/ConditionalAtoms.lean.
\begin{lemma}\label{lem:upper-smoothing}
For a genuine boundary pair with join $v$ and any legal continuation
$x$, there is a finite reveal containing the labeled prefix-$x$ rows
and all their continuations such that both interval endpoints are
known translates of the same hypergeometric variable. There is a
good event measurable in that reveal whose complement has probability
at most
\[
 3\exp\left\{-\frac{1-\log2}{2}g\min(L_v,R_v)\right\}.
\]
On this event the conditional variance is at least
\[
 \frac{g\min(L_v,R_v)}{16(2|v|+\chi^{-1})}.
\]
Thus every conditional endpoint atom has mass at most
$16\sqrt{(2|v|+\chi^{-1})/(g\min(L_v,R_v))}$.
\end{lemma}
\begin{proof}
Choose columns attaining $L_v$ and $R_v$. The same columns attain
these side means for every continuation $x$, since their
first-difference pieces are unchanged. It suffices to describe the
case $L_v\le R_v$; reflection interchanges the two sides, and the
complement of a hypergeometric sample has the same variance.

Suppose the left maximizing column is later than the right maximizing
column $j$. Expose every column except $j$. Let $T$ be the rows of
$R$ matching $x$ before $j$, and let $B\subseteq T$ be the rows whose
full suffix through the end of $x$ is lexicographically less than
the corresponding suffix of $x$. Expose column $j$ on
$A=T\setminus B$. Set $p=\sum_{d\le x_j}p_d^{(j)}$ and let $a$
be the number of exposed digits in $A$ at most $x_j$. The number
of such digits in $B$ is conditionally
$\operatorname{Hyp}(np-a,n-|A|,|B|)$. All other contributions to the
lower endpoint are known; the upper endpoint differs by the known
prefix count. Every prefix-$x$ row lies in $A$, so its continuation
is indeed known.

The set $B$ contains the left maximizing cylinder with column $j$
omitted; its mean is $L_v/p_{x_j}^{(j)}$. Require this count and
the remaining marked and unmarked counts to be at least half their
respective means. Lemma~\ref{lem:upper-atoms} bounds the union of
their failures by the asserted three lower-tail terms: each relevant
mean is at least $gL_v$. On this event, the marked and unmarked
fractions are at least $p/2$ and $(1-p)/2$, and the sample-complement
fraction is at least $g$, because $A\cup B\subseteq R$.
The variance is consequently at least
$g|B|p(1-p)/4\ge gL_vR_v/(8M)$.

If the left maximizing column $j$ is earlier, take $B$ to consist of
the rows with larger full suffix through the end of $x$,
and put $p=\sum_{d<x_j}p_d^{(j)}$. The same reveal leaves a
hypergeometric variable with marked count $np$ less the exposed
marked count. Now $B$ contains the right maximizing cylinder with
column $j$ omitted. The identical half-mean estimates again give
variance at least $gL_vR_v/(8M)$ and the same failure bound.

If the maximizing columns agree at $j$, expose all other columns and
the entire global set with digit $x_j$ in column $j$. Write
$p=\sum_{d<x_j}p_d^{(j)}$, $q=\sum_{d>x_j}p_d^{(j)}$, and let
$T$ be the rows of $R$ matching before $j$. On the unrevealed
$n(p+q)$ positions, exactly $np$ are marked. With $B$ the rows of
$T$ outside the exposed equal-digit set, the remaining variable is
$\operatorname{Hyp}(np,n(p+q),|B|)$.
Put $a=M\prod_{i<j}p_{x_i}^{(i)}$, so $L_v=ap$ and $R_v=aq$.
Require $|B|\ge a(p+q)/2$ and at least $(n-M)(p+q)/2$
non-equal digits outside $R$. These are two unconditional cylinder
counts with the indicated means. Each mean is at least $gL_v$,
so their failure probability is at most
$2\exp\{-(1-\log2)gL_v/2\}$. On this event the hypergeometric
mean is at least $L_v/2$, its unmarked fraction is at least $1/2$,
and its sample-complement fraction is at least $g/2$. Its variance
is at least $gL_v/8$.
In the unequal cases use \eqref{eq:upper-side-scale}; in the equal
case use $2|v|+\chi^{-1}\ge1$. The claimed common variance bound
follows. Applying the explicit atom estimate $4/\sigma$ proves the
last assertion.
\end{proof}

\subsection{Digit costs and stable intervals}
Fix $\delta>0$. Within each nonterminal boundary join, stop a branch
at its first prefix $x$ for which $\ell_x-F_v<\delta$; an immediate
stop is allowed. The terminal margin ensures that every branch stops:
at full depth, $\ell_x\le-\eta$ and
$\ell_x-F_v\le-\eta/2+\log(1/\chi)/(2\log N)<\delta$ for large $N$.
A proper stop has a parent with exponent at least $F_v+\delta$.
It also satisfies $F_v\ge\delta-\log(1/\chi)/\log N$.

For an extension from $v$ to $x$, let $S_x$ and $Q_x$ be its
normalized logarithmic costs on digits of mass at least $\chi$ and
less than $\chi$, respectively. Then
\begin{equation}\label{eq:upper-cost-split}
 \ell_x=V_v-S_x-Q_x,\qquad S_x,Q_x\ge0.
\end{equation}
Retain boundary words with at most
$(\eta+7)\log N/\log(1/\chi)$ exceptional positions before their
crossings, and retain extensions with at most
$(\log(1/\chi))^{-1/2}\log N$ exceptional positions.
Group the retained prefixes by their first pile, boundary words,
length, exceptional positions and digits, good-cost bin of width
$\eta/256$, and immediate/proper and large/small status. Here large
means $\ell_x\ge2\delta$. Terminal classes are singletons.

% Lean: FirstMoment/Count.lean, RetainedClasses.lean, Discard.lean,
% Entropy.lean, ClassExponents.lean, UniformClassDensity.lean.
\begin{lemma}\label{lem:upper-classes}
Take $\delta=\eta/256$. For sufficiently small fixed $\chi>0$ and
all sufficiently large $N$, the number of retained classes is at most
$N^{\eta/64}$ and the expected number of discarded rows in each
first pile is at most $N^{-(\eta+4)}$.
For a retained class $\mathcal C$, write
$H_{\mathcal C}=\log_N|\mathcal C|$. Every nonterminal class obeys
\begin{equation}\label{eq:upper-class-entropy}
 H_{\mathcal C}\le\theta S_x-B^{(N)}_\theta(M_v/\log N,M_x/\log N)
       +\frac{\eta}{64}
       +\frac{4\log k}{\sqrt{\log(1/\chi)}}
       \qquad(3\le\theta\le4).
\end{equation}
If $p_x(m)=\Pp((m,m+1)\in G_x)$, then every retained class satisfies
\begin{equation}\label{eq:upper-class-density}
 |\mathcal C|\sum_m p_x(m)^2\le N^{-\eta/2}
                  \qquad(x\in\mathcal C).
\end{equation}
All these estimates are uniform in the original array, first-column
assignment, and retained class.
\end{lemma}
\begin{proof}
A boundary is specified by its crossing position and digit and its
exceptional positions and digits; every other digit is the fixed
left or right default. For a word of length at most $T\log N$,
the number of ways to choose at most $q$ exceptional positions and
digits is at most $\sum_{j\le q}\binom{\lfloor T\log N\rfloor}j k^j$.
The bound $\binom mj\le(em/j)^j$ shows that the logarithm of this
quantity, divided by $\log N$, tends uniformly to zero as the
allowed exceptional fraction tends to zero. Lengths and crossing
times cost only powers of $\log N$, and the number of cost bins
is bounded for fixed $\chi$. Applying this to both boundaries and
the extension gives the class bound after choosing $\chi$ small.

For one row, independence of its digits gives probability at most
$(k\chi)^j$ for exceptional digits at $j$ prescribed positions.
A union bound, followed by the same binomial estimate, bounds the
expected discarded-row count. The boundary threshold contributes
logarithmic exponent at most $1-(\eta+7)+o_\chi(1)$; the extension
threshold contributes at most
$1-\sqrt{\log(1/\chi)}+o_\chi(1)$. Their sum is at most
$N^{-(\eta+4)}$ for sufficiently small $\chi$ and large $N$.

Within a class, the nonexceptional digits determine the word.
Their product tilted law, with digit probabilities proportional to
$(p_d^{(j)})^\theta$, assigns to $x$ the mass
$N^{-\theta S_x+B^G_\theta}$, where $B^G_\theta$ sums only over
nonexceptional positions. All good costs in the class lie in one
bin. Summing these tilted masses, whose total is at most one, yields
\eqref{eq:upper-class-entropy}: the bin costs at most $4(\eta/256)$,
and each removed coordinate costs at most $4\log k/\log N$.

Every internal edge can be charged to a distinct pair of equal rows.
The fixed-count falling-factorial bound therefore gives
\[
 \sum_m p_x(m)\le\tfrac12\mu_x^2N^{e_x},\qquad
 \sum_m p_x(m)^2\le\tfrac14 N^{4\ell_x+2e_x}.
\]
For a full word the latter sum is also at most $\mu_x$.
For a nonterminal block, condition as in
Lemma~\ref{lem:upper-smoothing}. Its internal edge offsets are known
under that reveal. For each offset, putting an edge at $m$ fixes
the lower endpoint. The atom bound thus bounds the maximal edge
probability on the good event by
$16\sqrt{(2|v|+\chi^{-1})/g}\,N^{-F_v}$ times the expected internal
edge count. Multiplying by the sum of edge probabilities, and
charging the bad event at most twice at each of the $n-1$ positions,
gives
\begin{align*}
 \sum_m p_x(m)^2
 &\le4\sqrt{(2|v|+\chi^{-1})/g}\,N^{4\ell_x-F_v+2e_x}\\
 &\quad+6(n-1)\exp\left\{-\frac{1-\log2}{2}g\min(L_v,R_v)\right\}.
\end{align*}

Put $q=\log(1/\chi)/\log N$ within this proof.
For every $\theta\in[3,4]$, the identity
\[
 4\ell_x-F_v+\theta S_x
 =\theta V_v+(4-\theta)(\ell_x-F_v)
                         +(3-\theta)F_v-\theta Q_x
\]
and $2F_v\ge V_v-q$ bound the right side by
$(\theta+3)\beta^{(N)}(M_v/\log N)/2+\delta+q/2$.
Using \eqref{eq:upper-class-entropy} and then taking the infimum
in \eqref{eq:upper-empirical-margin} bounds the large-block exponent
$H_{\mathcal C}+4\ell_x-F_v+2e_x$ by
\[
 -\eta+\delta+q/2+\eta/64+4\log k/\sqrt{\log(1/\chi)}.
\]
For a proper small block the parent stopping condition instead gives
$S_x\le V_v-F_v+q\le\beta^{(N)}(M_v/\log N)/2+3q/2$.
Its nonnegativity implies $\beta^{(N)}(M_v/\log N)\ge-3q$, whence
\[
 \theta S_x\le\tfrac{\theta+3}{2}\beta^{(N)}(M_v/\log N)+21q/2.
\]
Since $4\ell_x\le8\delta$, the small-block exponent is at most
$-\eta+8\delta+21q/2+\eta/64+4\log k/\sqrt{\log(1/\chi)}$.
For an immediate small block the class is a singleton and
$\ell_x+e_x\le\beta^{(N)}(K/\log N)\le-\eta$, because
$\psi_p(2)\ge a_p$. Its exponent is at most $-2\eta+4\delta$.
A terminal block contributes at most $N^{-\eta}$.

Choose $\chi$ also so that
$4\log k/\sqrt{\log(1/\chi)}\le\eta/256$, and then take $N$
large enough that $q\le\delta$. The large and proper-small
exponents above are at most $-3\eta/4$. The square-root prefactor
is at most $N^{\eta/16}$ eventually. In the large case, stability
and $\ell_x\ge2\delta$ imply $F_v>\delta$, so its failure term,
even multiplied by the total number $(K+1)k^K$ of possible prefixes,
is at most $N^{-\eta}$. Finally absorb the sum of the two terms
using $2\le N^{3\eta/16}$. This proves
\eqref{eq:upper-class-density} in all four cases.
\end{proof}

\subsection{Estimating the first moment}
% Lean: Localization/EndpointWindow.lean; Localization/RelevantParents.lean; Localization/Eventual.lean;
% FirstMoment/ActiveCount.lean, GraphAssembly.lean: actualGraphAssembly;
% FirstMoment/UniformGraph.lean: empiricalGraphDecay.
\begin{proof}[Proof of Proposition~\ref{prop:upper-first}]
We explain the localization needed to sum the class estimates.
Put $d=K-a-1$. Express the endpoints using the smaller side of a
genuine join and its cylinder count. Every one of these at most
$d+1$ terms has mean at most $\chi^{-1}N^{2F_v}$.
Lemma~\ref{lem:upper-atoms} and the union bound show that the
probability either endpoint differs from its mean by more than $w$
is at most
\[
 2(d+1)\exp\left\{-\frac{w^2}
                {4(d+1)^2(\chi^{-1}N^{2F_v}+w)}\right\}.
\]
For a proper stop, take $w=N^{F_v+\delta/4}$ and choose $N$ so that
$\log(1/\chi)/\log N\le\delta/2$. Then $F_v\ge\delta/2$,
and the last bound is at most $\exp\{-N^{\delta/8}\}$ uniformly.
Call a block relevant at $m$ when its interval contains $m$ with
probability at least $\exp\{-N^{\delta/16}\}$. This probability
exceeds the endpoint failure bound. Since the block lies in its
parent, some realization simultaneously places $m$ in the parent
and both parent endpoints within $w$ of their means. Thus $m$ lies
in the parent's mean interval enlarged by $w$ on each side.

The parent of a proper stop has mean length at least
$N^{F_v+\delta}>2w$. At a fixed depth the mean intervals are
ordered and have disjoint interiors. At most two such enlarged
intervals can contain $m$: if three did, the entire middle interval
would fit in an interval of length $2w$. Each parent has at most
$k$ children. Summing over depths gives at most $2kd$ relevant
proper blocks for each boundary descriptor, with at most one
additional immediate or terminal block. Thus at any edge the number of relevant retained
blocks is at most
$N^{\eta/64}(2k(K-a-1)+1)\le N^{\eta/32}$.
The union bound over all possible prefixes makes the total omitted
nonrelevant marginal at most $N^{-(\eta+4)}$ for large $N$.

For completeness, the finite summation behind this reduction is as
follows. Let $q_x(m)$ be the retained-block edge marginals and let
the remaining error marginal at $m$ be the discarded-row marginal
plus the total omitted marginal. At most $L$ retained marginals are
active at one edge. Cauchy--Schwarz and $(u+v)^2\le2u^2+2v^2$ give
\[
 \E|G\cap G'|
 \le4k^2 B+4(n-1)z^2+2L\sum_x\sum_m q_x(m)^2,
\]
where $B$ bounds the expected discarded-row count in every first
pile and $z$ bounds the omitted marginal. The factor $k^2$ follows
by summing the pile errors before squaring. This calculation uses
$\E|G\cap G'|=\sum_m\Pp(m\in G)^2$, which follows from independence.

Take $B=z=N^{-(\eta+4)}$ and $L\le N^{\eta/32}$.
Equation~\eqref{eq:upper-class-density}, summed over the at most
$N^{\eta/64}$ classes, bounds the double sum by
$N^{\eta/64-\eta/2}$. Hence, for sufficiently large $N$,
\[
 \E|G\cap G'|\le
 N^{-\eta}+2N^{3\eta/64-\eta/2}\le N^{-\eta/8}.
\]
All choices were made before $n,a,K$ and the column counts. Averaging
over the first-column assignments therefore preserves the estimate.
\end{proof}
% UPPER FIRST MOMENT END

% UPPER EXPONENTIAL BEGIN
\section{The exponential moment}\label{sec:upper-exponential}
\subsection{Preparatory estimates}
Fix $k\ge2$, $T>0$, $0<\eta\le1$, $0<g<1/k$ and
$0<\gamma\le1$. We consider a retained array of length at most
$T\log N$, with empirical margin at least $2\eta$. Its first column
is $g$-good, two further good columns occur within the next
$D=\lceil6/\gamma\rceil$ columns, and every consecutive window of
$\lceil\rho\log N\rceil+1$ columns contains at least
$\gamma\rho\log N-2$ good columns. All randomness below is in the
uniform assignments of these fixed column counts.

The short tail prefix has length $\min\{D,K-1\}$. The two good
columns in this prefix provide three
ordered short words $v_-<v_0<v_+$, each of tail mass at least
$g^2k^{-D}$. To see this, choose two ordered digits of mass at least
$g$ in each of these columns and a maximal-mass digit in every other
coordinate. Three of the four resulting words are strictly ordered
and have the required masses. Choose the lexicographically first
admissible triple; it depends only on the counts. Assign a full word
to the forward direction when its short tail prefix is at most $v_0$,
and to the backward direction when it is at least $v_0$.

In a first-column pile $p$ of size $n_p$, reveal the sorted words
from the left for the forward direction and from the right for the
backward direction. The natural history $h_i$ consists precisely of
the first $i$ words so revealed. The chronological order of their
coordinates is unchanged in the backward exploration. We never add
a global concentration event to a natural history.

Let $H$ be the following event on one matrix. For every pile and
later column, each digit count outside the pile differs from its
mean by at most $N^{2/3}$. For every prefix $x$ with
$\mu_x\ge N^{\eta/10}$, its count differs from $\mu_x$ by at most
$\mu_x^{2/3}$, and each later digit count within that prefix differs
from its count times the original digit proportion by at most
$\mu_x^{2/3}$. Finally, in every pile of size at least $N^{\eta/10}$,
the numbers of words strictly before and at or before $v_0$ differ
from their means by at most $g^2k^{-D}n_p/4$.

% Lean: Exploration/UniformConcentration.lean: uniformConcentration;
% Exploration/HistoryProbability.lean: historyGoodProbability;
% Exploration/Regularity.lean: historyGood_unionBound.
\begin{lemma}\label{lem:upper-history-good}
There exist $A,c_H>0$ and $N_0$, depending only on the fixed
parameters, such that for $N\ge N_0$,
\begin{equation}\label{eq:upper-H}
 \Pp(H^c)\le AN^A\exp\{-c_H N^{\eta/30}\}.
\end{equation}
For any fixed pile, direction and reveal time, the probability of
histories $h_i$ for which
\begin{equation}\label{eq:upper-history-good}
 \Pp(H^c\mid h_i)\le
          \sqrt{AN^A\exp\{-c_H N^{\eta/30}\}}
\end{equation}
fails is at most the right side of \eqref{eq:upper-history-good}.
The probability that it fails at any nonterminal positive reveal
time in any pile or direction is at most $2N$ times that bound.
\end{lemma}
\begin{proof}
Use the stronger simultaneous event obtained by prescribing an
arbitrary set of allowed digits in every column. If such a cylinder
has mean $\mu$, require its count to lie within
$\max\{\mu,N^{\eta/10}\}^{2/3}/2$ of its mean.
The transform proof of Lemma~\ref{lem:upper-atoms} applies equally
to allowed-digit sets. Substituting this tolerance in its tail
bound shows that each test fails with probability at most
$2\exp\{-N^{\eta/30}/24\}$.
There are exactly $(2^k)^K$ tests. Their simultaneous failure
probability is therefore at most
$2(2^k)^K\exp\{-N^{\eta/30}/24\}$.

This stronger event implies $H$ deterministically. Outside-pile
digit counts are allowed-digit cylinders of mean at most $N$.
For a prefix of mean $\mu_x\ge N^{\eta/10}$, both its count and
the count obtained by prescribing one later digit have errors at
most $\mu_x^{2/3}/2$. Subtract the original digit proportion times
the first error from the second; their sum is at most
$\mu_x^{2/3}$. Each guard is a sum of at most $k^D$ short-cylinder
counts. In a pile with $n_p\ge N^{\eta/10}$, all their means are
at most $n_p$, so its error is at most $k^D n_p^{2/3}/2$.
For large $N$,
$2k^D/(g^2k^{-D})\le N^{\eta/30}\le n_p^{1/3}$,
which makes this at most the required $g^2k^{-D}n_p/4$.
Finally $(2^k)^K\le N^{kT\log2}$. This proves
\eqref{eq:upper-H}, with the permissible explicit choices
$A=2+kT\log2$ and $c_H=1/24$.

The expectation of $\Pp(H^c\mid h_i)$ is $\Pp(H^c)$, by summing
over the finite history fibers. Markov's inequality proves the
single-history assertion. There are $2\sum_p n_p=2N$ pile-direction
reveal indices, so the final assertion follows by the union bound.
\end{proof}

Choose $0<\tau<g$ and write
$\alpha=\log(1/g)/\log(1/\tau)$. We take $\tau$ so small that
\begin{equation}\label{eq:upper-alpha}
 \alpha<\tfrac12,\qquad
 \alpha\le\frac{\eta/10}{4(1+\eta/10)},\qquad
 \alpha\le\frac{\eta}{40(1+\eta)}.
\end{equation}
This choice precedes $\rho,N$, the array, and the history. Denote by
$\lambda_x$ the product probability of a prefix $x$, so
$\mu_x=N\lambda_x$.

% Lean: Reservoir/PrefixCount.lean: prefixCount; Reservoir/ExtraCopies.lean: extraCurrentCopies;
% Reservoir/TailCollisions.lean: tailCollisions; Reservoir/BackwardCounts.lean.
\begin{lemma}\label{lem:upper-reservoir}
There exist $C>0$ and $N_0$, depending only on the fixed parameters,
with the following property for $N\ge N_0$.
Fix a positive-probability natural history ending at word $s$ in its
assigned direction, and suppose $\Pp(H^c\mid h_i)\le\varepsilon$,
where $\varepsilon\ge0$. Every prefix cylinder $x$ strictly beyond $s$ in this
direction and contained in its first pile satisfies
\begin{align}
 \E[C_x\mid h_i]&\le CN\lambda_x^{1-\alpha}+N\varepsilon,
                                           \label{eq:upper-reservoir-count}\\
 \E\left[\sum_{w\succeq x}\binom{C_w}{2}\,\middle|\,h_i\right]
 &\le CN^2\lambda_x^{2(1-\alpha)}
       \prod_{j>|x|}\sum_d(p_d^{(j)})^{2(1-\alpha)}
                  +N^2\varepsilon.          \label{eq:upper-reservoir-pairs}
\end{align}
If $t$ copies of the current word $s$ have been observed, then also
\begin{equation}\label{eq:upper-extra-copies}
 \E[(C_s-t)_+\mid h_i]
       \le CN\lambda_s^{1-\alpha}+N\varepsilon.
\end{equation}
\end{lemma}
\begin{proof}
We give the forward proof. Mark one or two labeled rows of the
initial pile and expose all other pile rows. Keep all outside-pile
rows unexposed. In each later column, if $m\in\{1,2\}$ rows are
marked, the remaining population has size $N-n_p+m$. Conditional
on the exposed values, its assignments are uniform with the exact
remaining counts and independent between columns. This follows
by counting completions of each column; the number of completions
is the corresponding multinomial coefficient.

Deleting marked rows which lie strictly after the observed prefix
leaves that prefix valid. Completing the marked rows inside the
right reserve $v_+$ therefore preserves the history. On every
exposed fiber meeting $H$, its residual reserve counts differ from
their means by at most $N^{2/3}+m$. Since $N-n_p\ge gN$ and
the reserve depth is fixed, the product of the relevant falling-
factorial ratios is at least $(g^2k^{-D})^m/2$ for large $N$.
This gives a positive, uniform lower bound for the history denominator
on such a fiber.

For the numerator, a prescribed digit of mass $p\ge\tau$ has
residual proportion at most $p(1+CN^{-1/3})$, by the outside-pile
concentration bound. A digit of mass $p<\tau$ has residual
proportion at most $p/g$, because the remaining population has
size at least $gN$. The identity $\tau^\alpha=g$ gives
$p/g\le p^{1-\alpha}$ in the latter case; in the former,
$p\le p^{1-\alpha}$. For $m=1,2$, the exact sampling factor is
no larger than the $m$th power of that residual proportion.
Multiplying over at most $T\log N$ columns costs at most $2^m$
for large $N$. The numerator is consequently bounded by
$2^m$ times the prescribed product mass to power $m(1-\alpha)$.

Divide by the reserve lower bound and then average the exposed
information back down to the natural history. Summing the one-row
bound over labeled rows gives \eqref{eq:upper-reservoir-count}.
Summing the two-row bound over pairs and their common full tail
gives \eqref{eq:upper-reservoir-pairs}; the full-tail sum factorizes
into the displayed product. The bad part is charged once, before
this summation: a count is at most $N$ and the total number of
pairs at most $N^2$, so its costs are $N\varepsilon$ and
$N^2\varepsilon$, respectively. In particular there is no error
factor counting terminal words.

For the current word, mark only a removable copy: after deleting it,
the $t$ observed copies must remain. The number of removable copies
dominates $(C_s-t)_+$. Complete it in the right reserve as before,
and apply the same one-row calculation. This proves
\eqref{eq:upper-extra-copies}, including the possible edge from the
last observed copy to the first unobserved copy. For the backward
direction use the left reserve and the reversed spatial inequalities.
The columns themselves retain their chronological order.
\end{proof}

% Lean: Reservoir/ConditionalFutureCylinderLaw.lean: conditionalFutureCylinderLaw,
% conditionalFutureCylinderLaw_unlabeled; Conditional/Jensen.lean;
% Conditional/GraphSquares.lean: fineBlockSquares_le_firstMoment.
\begin{lemma}\label{lem:upper-induced-law}
Let a finite family of disjoint prefix cylinders lie strictly beyond
the last word of a natural history. Refine the history by revealing
their labeled row sets, their rank intervals, and their counts of
each digit in each later column, together with the matrix outside
their interiors. Conditional on this information, the internal
graph of any cylinder has the law obtained from independent uniform
fixed-count tail columns with its induced counts, translated by its
known lower rank. The same statement holds after averaging out the
labeled row sets while retaining the induced counts and intervals.
\end{lemma}
\begin{proof}
Inside one specified cylinder all its prefix coordinates are fixed.
The natural history concerns words strictly outside it, so it places
no further order condition on its interior. After the outside matrix
and inside column counts are fixed, every compatible assignment of
one tail column inside has the same number of full completions.
Counting independently in each column therefore gives the uniform
product of its fixed-count column laws. Sorting inside the cylinder
sorts precisely these tail words; its lower rank supplies the
deterministic translation. Permuting the labels within the fixed
row set does not change that graph law. Hence every labeled atom
with the same counts and interval has the same conditional graph
law, and averaging these atoms proves the unlabeled assertion.
\end{proof}

\subsection{The exploration}
For independent matrices, let $X_i^{p,\pm}$ count the common
directional edges in pile $p$ with both endpoints among its first
$i$ revealed vertices. Use the identically zero process if
$n_p<N^{\eta/10}$. Each process starts at zero and has increments
zero or one. Its pair history consists of the two natural histories
in the same pile and direction. The remaining count is
$X_{n_p}^{p,\pm}-X_i^{p,\pm}$; at time $i$ it includes the edge
crossing from the last revealed vertex to the next vertex.
A pair history is good if each of its two histories satisfies
\eqref{eq:upper-history-good}; declare time zero and terminal time
good without further restrictions.

% Lean: Contracts/V1/Conditional.lean: C; Conditional/Main.lean: remainingCommonEdges;
% Conditional/OneHistory.lean: oneHistory_uniform; Conditional/Parameters.lean.
\begin{proposition}\label{prop:upper-conditional}
Under the fixed parameters and retained-array hypotheses of this
section, there is $\rho_0>0$ such that every fixed
$0<\rho\le\rho_0$ admits $A,c_H,C>0$ and $N_0\ge2$ for which
\eqref{eq:upper-H} holds and, on every positive-probability good
pair-history fiber in a pile of size at least $N^{\eta/10}$,
\begin{equation}\label{eq:upper-conditional}
 \E[X_{n_p}^{p,\pm}-X_i^{p,\pm}\mid h_i,h'_i]
                          \le CN^{-3\eta/80}
                         \qquad(0\le i\le n_p).
\end{equation}
At $i=n_p$ the remaining count is identically zero. All constants
are chosen before the array, pile, direction, and history.
\end{proposition}
\begin{proof}
Choose $\tau$ according to \eqref{eq:upper-alpha}, and take
$\rho\log k\le3\eta/160$. Fix one positive history with current
word $s$. If $s$ is outside its assigned direction, no directional
edge remains. Otherwise partition the strict future in that pile
into prefix cylinders. Following the two boundary paths of the
lexicographic interval gives at most $2kK$ disjoint cylinders.
If at least $\lceil\rho\log N\rceil+1$ columns remain after a
prefix, refine it to just before a good column in the next such
window. The integer-window condition and the three-index argument
used in Lemma~\ref{lem:upper-restart} provide that column. Each
refinement adds at most $\lceil\rho\log N\rceil+1$ coordinates.
The resulting disjoint family $\mathcal B$ therefore has size at
most
\begin{equation}\label{eq:upper-future-count}
 2k^3T(\log N)N^{\rho\log k}.
\end{equation}
Each member either has a good next column or has fewer than that
many remaining columns. In the latter case its maximal tail mass
is at least $k^{-\lceil\rho\log N\rceil-1}$, whereas its mean
times that tail mass is at most $N^{-\eta}$. Its mean is thus
less than $N^{-\eta/10}$ for large $N$.

Let $q_e(h_i)$ be the conditional probability of a remaining
directional edge at position $e$. Refine the history only for this
estimate, using the family information in
Lemma~\ref{lem:upper-induced-law}. Conditional Jensen gives
\[
 \sum_e q_e(h_i)^2
 \le\E\left[\sum_e q_e(h_i,Z)^2\,\middle|\,h_i\right].
\]
The refined edge slots split into the remaining current-word run
and the disjoint cylinder interiors. An equal-word edge cannot
cross between two distinct cylinders. The current-run contribution
is bounded by \eqref{eq:upper-extra-copies}. For the final cylinders,
after refinement, use three regimes.

If $\mu_x\le N^{-\eta/10}$, the sum of squared edge probabilities
is at most the expected vertex count. By
\eqref{eq:upper-reservoir-count} it is at most
\[
 CN^{-\eta/10+\alpha(1+\eta/10)}+N\varepsilon
                   \le CN^{-3\eta/40}+N\varepsilon.
\]
If $N^{-\eta/10}<\mu_x\le N^{\eta/10}$, it is at most the
expected number of equal-word pairs. Write
$q_* =\prod_{j>|x|}\max_d p_d^{(j)}$.
The terminal margin gives $\mu_xq_*\le N^{-\eta}$, hence
$q_*\le N^{-\eta+\eta/10}$. Since $2(1-\alpha)>1$,
\[
 \prod_{j>|x|}\sum_d(p_d^{(j)})^{2(1-\alpha)}
                         \le q_*^{1-2\alpha}.
\]
Equation~\eqref{eq:upper-reservoir-pairs} bounds this regime by
\[
 CN^{-\eta+3\eta/10+2\alpha(1+\eta-2\eta/10)}+N^2\varepsilon
                         \le CN^{-3\eta/5}+N^2\varepsilon.
\]

Finally, suppose $\mu_x>N^{\eta/10}$. Such a cylinder has a good
next column. On $H$, its size $m$ and later digit counts satisfy
\[
 \tfrac12\mu_x\le m\le2\mu_x,\qquad
 \max_{j,d}|p_{d,\mathrm{ind}}^{(j)}-p_d^{(j)}|
                            \le2N^{-\eta/30}.
\]
For two simplex vectors at sup-distance at most $z$, the scalar
cost bounds are
$|a_p-a_q|\le kz$ and
$|\psi_p(\theta)-\psi_q(\theta)|\le4k^4z$ for $2\le\theta\le4$.
Indeed, their maxima are at least $1/k$, their power sums are at
least $k^{-3}$, the logarithm is Lipschitz above these lower bounds,
and $|x^\theta-y^\theta|\le4|x-y|$ on $[0,1]$.
Thus replacing the suffix population by the induced one loses at
most
\[
 \frac72\left(\frac{\log2}{\log N}+2kT N^{-\eta/30}\right)
                         +24k^4T N^{-\eta/30}
\]
in its empirical expression. The maximal-prefix inequality
$\mu_x\le N\prod_{j\le|x|}\max_d p_d^{(j)}$ supplies the required
bound on its initial population exponent. The original margin
$2\eta$ therefore leaves margin $\eta$ in the induced array, and
its first column is $g/2$-good, for large $N$.
Lemma~\ref{lem:upper-induced-law} and
Proposition~\ref{prop:upper-first} bound its conditional
self-intersection by $N^{-\eta/8}$ on these good family atoms.
Their complement has conditional probability at most $\varepsilon$,
since the natural history is good and $H$ implies the needed
conditions. Its contribution is at most $N\varepsilon$.

The current word has mean at most $N^{-\eta}$, so
\eqref{eq:upper-extra-copies} bounds its excess by
$CN^{-\eta+\alpha(1+\eta)}+N\varepsilon$.
Combining the regimes, using \eqref{eq:upper-future-count}, and
taking $\varepsilon=\sqrt{AN^A e^{-c_HN^{\eta/30}}}$, gives
\[
 \sum_e q_e(h_i)^2\le
 CN^{-3\eta/40}
 +C(\log N)N^{\rho\log k}
       \bigl(N^{-3\eta/40}+N^{-3\eta/5}+N^{-\eta/8}
                                      +N^2\varepsilon\bigr)
 \le CN^{-3\eta/80}.
\]
The last step uses $\rho\log k\le3\eta/160$ and
$\log N\le N^{3\eta/160}$ eventually; the stretched exponential
absorbs every polynomial factor.

The two matrices remain independent when their separate histories
are fixed. Thus the expected remaining common count is
$\sum_e q_e(h_i)q_e(h'_i)$, bounded by half the sum of their two
self-intersections. This proves \eqref{eq:upper-conditional} at
positive nonterminal times. At time zero the remaining directional
count is bounded by the total common-edge count, so
Proposition~\ref{prop:upper-first} applies directly. The terminal
count is zero by definition. The backward proof uses the same
argument with the spatial order reversed.
\end{proof}

\subsection{Proof of the exponential estimate}
% Lean: Exponential/Stopping.lean; Exponential/Moments.lean;
% Exponential/Excess.lean: adaptedExcess.
\begin{lemma}\label{lem:upper-adapted}
On a finite probability space let $(\mathcal F_i)_{0\le i\le m}$
be an increasing filtration, let $X_0=0$ be an adapted process with
increments in $\{0,1\}$, and let $B_i\in\mathcal F_i$.
Suppose that for every $i<m$, on every positive-probability
$\mathcal F_i$-atom contained in $B_i$,
\[
 \E[X_m-X_i\mid\mathcal F_i]\le u,
 \qquad u\ge0.
\]
For $t\ge0$ with $(e^t-1)u<1$ and any event
$C\subseteq\bigcap_{i=0}^m B_i$,
\begin{equation}\label{eq:upper-adapted}
 \E[(e^{tX_m}-1)\mathbf1_C]
                  \le\frac{(e^t-1)u}{1-(e^t-1)u}.
\end{equation}
\end{lemma}
\begin{proof}
Set $D_i=\bigcap_{h=0}^iB_h$ and
$Z_i=\sum_{j=1}^i(X_j-X_{j-1})\mathbf1_{D_j}$.
Then $Z$ is adapted, has unit increments, and
$Z_m-Z_i\le\mathbf1_{D_i}(X_m-X_i)$. Therefore its conditional
expected remainder is at most $u$ on every positive-probability
history atom, without a good-event restriction.

For every integer $r\ge0$ we claim
$\E[\binom{Z_m-Z_i}{r}\mid\mathcal F_i]\le u^r$.
The claim is immediate for $r=0$. For $r\ge1$, the pathwise identity
\[
 \binom{Z_m-Z_i}{r}=
 \sum_{j=i+1}^m(Z_j-Z_{j-1})\binom{Z_m-Z_j}{r-1}
\]
counts an $r$-tuple of future successes by its first success.
Condition each summand on $\mathcal F_j$, apply the induction
hypothesis to the second factor, then condition on $\mathcal F_i$.
The conditional first-moment bound proves the claim.
The binomial theorem now gives
\[
 \E(e^{tZ_m}-1)
 \le\sum_{r\ge1}((e^t-1)u)^r
       =\frac{(e^t-1)u}{1-(e^t-1)u}.
\]
On $C$ no increment was discarded, so $Z_m=X_m$; outside $C$
the exponential excess is nonnegative. This proves the result.
\end{proof}

% Lean: Exponential/FiniteTV.lean: finiteExplorationTV_of_majorization,
% finiteExplorationTV; Exponential/Aggregation.lean: finiteExplorationExcess.
\begin{proposition}\label{prop:upper-exponential}
Let $S$ be an event on one matrix whose graph components have size
at most $B\ge2$. Suppose $d\ge1$ separate explorations on the pair
of matrices satisfy Lemma~\ref{lem:upper-adapted} with common bound
$u$, and that on $S\times S$ all their good events hold and
$|G\cap G'|\le\sum_{a=1}^d X_{m_a}^{(a)}$.
Put $b=\binom B{\lfloor B/2\rfloor}$. If $(b^d-1)u<1$, then
\[
 \dTV(P_N^K,\nuUnif)
 \le\frac12\sqrt{\frac{(b^d-1)u}{1-(b^d-1)u}}+\Pp(S^c).
\]
\end{proposition}
\begin{proof}
For nonnegative $x_1,\ldots,x_d$, convexity of the exponential gives
$b^{\sum_a x_a}-1\le d^{-1}\sum_a(b^{d x_a}-1)$.
Apply this on $S\times S$ and then apply
Lemma~\ref{lem:upper-adapted} separately to each exploration with
$t=d\log b$ and $C=S\times S$. It bounds the expected truncated
excess $b^{|G\cap G'|}-1$ by $(b^d-1)u/(1-(b^d-1)u)$.
Lemmas~\ref{lem:upper-overlap} and \ref{lem:upper-truncation}
then give the asserted total variation bound. No common filtration
for the separate explorations is required.
\end{proof}

% Lean: Exploration/UniformRegularity.lean, DirectionalCover.lean, SmallPiles.lean;
% EmpiricalUpper/ExplorationTV.lean: concreteExplorationTV.
We finish by checking the truncation used in the proof of
Theorem~\ref{thm:upper-empirical}. Let $\mathcal R$ be the event
that $H$ holds, every positive nonterminal history in every pile
and direction satisfies \eqref{eq:upper-history-good}, and every
pile smaller than $N^{\eta/10}$ is edge-free. The expected number
of equal-word pairs in such a pile is at most
$n_p\mu_*/2$. Summing over at most $k$ small piles gives
\[
 \Pp(\mathcal R^c)\le AN^A e^{-c_HN^{\eta/30}}
      +2N\sqrt{AN^A e^{-c_HN^{\eta/30}}}
      +\frac{k}{2}N^{\eta/10-2\eta}.
\]
The guard concentration ensures the directional cover for two
matrices: their middle-cylinder rank intervals overlap, because
their common expected gap is at least $g^2k^{-D}n_p$ whereas the
two endpoint errors total at most $g^2k^{-D}n_p/2$.
Thus an edge common to both matrices is assigned to a common
direction in its pile. Small piles contribute no edges on
$\mathcal R$. Hence
$|G\cap G'|\le\sum_{p,\pm}X_{n_p}^{p,\pm}$ on
$\mathcal R\times\mathcal R$.
The natural pair histories refine with time, the edge processes
are adapted and have unit increments, and their good events are
history measurable. Proposition~\ref{prop:upper-conditional}
provides every conditional remainder bound, including time zero.
Intersecting $\mathcal R$ with the bounded-component event therefore
verifies all hypotheses of Proposition~\ref{prop:upper-exponential}.
This completes the upper-bound argument.
% UPPER EXPONENTIAL END

